% Phase 3A: one-factor synthetic robustness (profile robustness_phase3).
% Figures: representative RMSE vs rho curves at baseline-like knob values.
% Main table: robustness_tables_fragment_measurement_ratio.tex.
% Appendix: robustness_tables_fragment_rho060.tex (regenerate via build_robustness_tables.py).

\section{Robustness: synthetic one-factor ablations (Phase~3A)}
\label{sec:exp_robustness}

We stress the hybrid estimators along single axes of the synthetic generator while keeping all other Phase~0 settings fixed.
Each row in the appendix tables aggregates ten seeds; the main text shows a compact \emph{curve} view at representative knob values and the full $\rho$ sweep for the measurement-ratio axis (where each run fixes one $\rho$).

\paragraph{Protocol.}
Runs are produced by \texttt{python sir\_cs\_pipeline\_optimized.py --profile robustness\_phase3} with \texttt{--robustness-axis} and \texttt{--robustness-value}; the sweep manifest is \texttt{outputs/robustness\_phase3/sweep\_manifest.csv}.

\begin{figure}[t]
  \centering
  \begin{minipage}[t]{0.48\linewidth}
    \centering
    \includegraphics[width=\linewidth]{figures/robustness_phase3/residual_k/v_6/runs/20260418_190205/01_rmse_vs_measurement_ratio.png}\\[-0.5em]
    {\small (a) Varying $k$ at $k=6$ (baseline match).}
  \end{minipage}\hfill
  \begin{minipage}[t]{0.48\linewidth}
    \centering
    \includegraphics[width=\linewidth]{figures/robustness_phase3/measurement_noise_std/v_0p02/runs/20260419_082615/01_rmse_vs_measurement_ratio.png}\\[-0.5em]
    {\small (b) Varying $\sigma_\eta$ at $\sigma_\eta=0.02$ (baseline).}
  \end{minipage}\\[0.8em]
  \begin{minipage}[t]{0.48\linewidth}
    \centering
    \includegraphics[width=\linewidth]{figures/robustness_phase3/residual_amplitude/v_1p2/runs/20260419_210918/01_rmse_vs_measurement_ratio.png}\\[-0.5em]
    {\small (c) Varying innovation amplitude at $1.2$ (baseline).}
  \end{minipage}\hfill
  \begin{minipage}[t]{0.48\linewidth}
    \centering
    \includegraphics[width=\linewidth]{figures/robustness_phase3/output_noise_std/v_0p01/runs/20260420_083135/01_rmse_vs_measurement_ratio.png}\\[-0.5em]
    {\small (d) Varying output noise std at $0.01$ (baseline).}
  \end{minipage}
  \caption{Representative RMSE versus $\rho=m/N$ curves under one-factor robustness sweeps. Methods include classical and LFISTA hybrids plus baselines from the same protocol.}
  \label{fig:robust_rmse_panel}
\end{figure}

Figure~\ref{fig:robust_rmse_panel} supports the same ordering seen in the dedicated LFISTA comparison: LFISTA hybrids track below \texttt{hybrid\_fista} across $\rho$ in these representative regimes.

% Auto-generated by paper/build_robustness_tables.py (do not edit by hand).

\begin{table}[t]
  \centering
  \caption{Test RMSE when the sweep fixes a single $\rho=m/N$ (one column per run). At $\rho=0.80$, \texttt{hybrid\_fista} is marginally better than \texttt{hybrid\_lfista\_joint}.}
  \label{tab:robust_measurementratio}
  \small
  \begin{tabular}{@{}lcccc@{}}
    \toprule
    $\rho=m/N$ & \texttt{hybrid\_fista} & \texttt{hybrid\_lfista\_joint} & \texttt{hybrid\_lfista\_frozen} & \texttt{ml\_only} \\
    \midrule
    0.10 & 0.353 & 0.335 & 0.335 & 0.369 \\
    0.20 & 0.336 & 0.320 & 0.320 & 0.369 \\
    0.30 & 0.319 & 0.300 & 0.300 & 0.369 \\
    0.40 & 0.296 & 0.277 & 0.277 & 0.369 \\
    0.50 & 0.272 & 0.254 & 0.255 & 0.369 \\
    0.60 & 0.245 & 0.229 & 0.231 & 0.369 \\
    0.80 & 0.180 & 0.181 & 0.186 & 0.369 \\
    \bottomrule
  \end{tabular}
\end{table}



\paragraph{Discussion.}
The measurement-ratio sweep in Table~\ref{tab:robust_measurementratio} isolates very high sampling ratios; at $\rho=0.80$ the classical FISTA hybrid slightly improves over the joint LFISTA model, whereas LFISTA remains clearly better than \texttt{ml\_only}.
Appendix~\ref{sec:appendix_robust_tables} reports analogous RMSE summaries at $\rho=0.60$ for $k$, measurement noise, innovation amplitude, and output noise.
